\begin{textsample}{POS Dim 1 – pi – Score 60.00 – pi/pi\_em\_11.txt}  \label{ex:f1_pos_081}
3.1.2 Ensino e aprendizagem O processo de ensino e aprendizagem tem sido objeto de reflexão há décadas.
A preocupação com o ato de \textbf{educar} vem desde a antiguidade.
O filósofo grego Sócrates, nos anos entre 469 a 399 a.C., opondo -se à filosofia \textbf{que} procurava explicar o mundo \textbf{baseado} na observação da natureza, defendeu a proposição de \textbf{que} o ser humano deveria conhecer a si mesmo, tendo como inquietação, levar as pessoas, por meio do autoconhecimento, à sabedoria e à prática do bem.
Para ele a valorização do diálogo era algo importantíssimo, pois serviria como meio de investigação acreditando, também, \textbf{que} a aprendizagem deveria começar a partir da educação do corpo, \textbf{que} por sua vez, permitiria controlar o físico.
Por volta dos anos 427 a 347 a.C., a formação do \textbf{homem} moral, como objetivo final da educação, assim como a defesa de \textbf{um} estado \textbf{justo} para acolher cada indivíduo foram objetivos defendidos por Platão.
Isso fez com \textbf{que} sua \textbf{pedagogia} se \textbf{aproximasse} da sua própria filosofia, \textbf{dizendo} \textbf{que} a busca da verdade é mais importante \textbf{que} a dos dogmas incontestáveis.
Este pensador defendeu a ideia de \textbf{que} professor ou aluno pensem sobre si próprios.
Aristóteles ( 384 a 322 a.C. ), discípulo de Platão, acreditava \textbf{que} a educação deveria ser tomada como o real caminho para a vida pública, tendo em vista \textbf{que} somente com \textbf{uma} adequada educação seria possível formar o caráter do aluno.
Neste sentido, este filósofo defendia a ideia de \textbf{que} a criança aprende com o ato de imitar os bons hábitos, e consequentemente, se tornar \textbf{um} cidadão ético e voltado à efetividade do bem comum em seu meio social.
No período medieval, com o paradigma teocêntrico em alta nos mais diversos segmentos ( religioso, cultural, político, epistemológico ), a educação era tratada como especificidade da Igreja, assim como a expansão da mesma na Europa.
Era responsabilidade da Igreja a disseminação da cultura, bem como dos processos educacionais deste contexto histórico, destacando -se \textbf{um} prévio curriculum instrucional \textbf{embasado} nas sete artes liberais, compostas por : gramática, retórica, lógica, aritmética, geografia, astronomia e música.
A Igreja torna -se a detentora do conhecimento e de formas educacionais deste período, tendo em vista \textbf{que} somente esta instituição poderia dar aval, conforme seus interesses, do \textbf{que} poderia ser ou não lecionado.
O ensino era pautado na leitura de textos e \textbf{exposição} das principais ideias pelo professor, formando muitas vezes debates entre professores e alunos ( scholastica disputattio ), e \textbf{que} se tornou \textbf{conhecido} como escolástica, \textbf{método} de estudos muito usado por São Tomás de Aquino ( 1225-1274 ).
No início da modernidade, é possível ser observada \textbf{uma} importante transição no \textbf{que} se refere aos \textbf{métodos} educacionais, onde o \textbf{método} escolástico ( medieval ) ganha \textbf{uma} remodelagem fundamentando -se não mais no conhecimento \textbf{revelado} ( religioso ), mas na \textbf{vertente} \textbf{que} defende \textbf{que} o próprio ser humano, \textbf{enquanto} ser racional, tem aparato mental e pode estabelecer critérios específicos para o alcance do conhecimento verdadeiro.
Este intermediado pelas experiências e uso coerente da razão, não mais tendo a razão como \textbf{mero} instrumento da religião.
Neste sentido, podem -se destacar algumas perspectivas educacionais \textbf{que} têm sua origem e motivações a partir de significativos acontecimentos e mudanças ocasionados em determinados momentos da história \textbf{moderna}, \textbf{influenciando} os mais diversos aspectos \textbf{antropológicos}, tais como econômico, político, social e cultural.
Com isto, pode ser afirmado \textbf{que} na modernidade têm -se ênfases de propostas educacionais direcionadas a \textbf{métodos} de ensino pensados a partir de \textbf{fundamentos} filosóficos, conforme se observa na continuidade deste texto.
Em \textbf{uma} perspectiva educacional pós contexto da Revolução Francesa – na qual se defendeu \textbf{uma} sociedade movida pelos princípios da Liberdade, Igualdade e Fraternidade, como direito a ser usufruído por toda sua população, surge \textbf{uma} marcante discussão a partir das ideias de Karl Marx ( 1818-1883 ), onde se tem a defesa de \textbf{um} ensino pautado na transformação social.
Na visão deste \textbf{autor} a sociedade pode ser melhor compreendida como \textbf{uma} composição de relações sociais, onde o ser humano é feito a partir de suas relações com a alteridade, ou seja, com os outros \textbf{homens}.
Assim, o ser humano seria movido à base de suas necessidades e liberdade, existindo \textbf{uma} \textbf{dialética} em \textbf{que} se possam justificar os fins.
No \textbf{século} XX, teve grande contribuição na formação da visão pedagógica, os pensamentos de Maria Montessori ( 1870-1952 ), ao defender a ideia de \textbf{que} a criança conquistava a educação, partindo da concepção de \textbf{que} já nascemos com a capacidade de ensinar a nós mesmos, desde \textbf{que} nos sejam dadas as condições.
Seus estudos \textbf{basearam} -se na observação de \textbf{que} as crianças aprendem melhor pela experiência direta da procura e descoberta, cabendo ao professor acompanhar o processo e detectar o modo particular de cada \textbf{um} manifestar seu potencial.
Ao longo do \textbf{século} XX, diferentes concepções pedagógicas fundamentaram o processo de ensino e aprendizagem, tais concepções, segundo Saviani ( 2006 ), podem ser agrupadas em duas grandes tendências : a primeira seria composta pelas concepções pedagógicas \textbf{que} dariam prioridade à \textbf{teoria} sobre a prática e a segunda tendência, inversamente, compõe -se das concepções \textbf{que} subordinam a \textbf{teoria} à prática e dissolvem a \textbf{teoria} na prática.
As modalidades de \textbf{pedagogia} tradicional, de \textbf{vertente} religiosa ou leiga, dão prioridade à \textbf{teoria} sobre a prática, e \textbf{focam} especificamente nas \textbf{teorias} de ensino.
Ao passo \textbf{que} as modalidades de \textbf{pedagogias} novas priorizam a prática, cuja ênfase é posta nas “ \textbf{teorias} da aprendizagem ”.
Ainda conforme Saviani ( 2006 ), as concepções tradicionais, desde a \textbf{pedagogia} de Platão e a \textbf{pedagogia} cristã, passando pelas \textbf{pedagogias} dos humanistas, assim como a \textbf{pedagogia} idealista de Kant, Fichte e Hegel, o humanismo racionalista, a \textbf{teoria} da evolução pautavam -se pela \textbf{centralidade} da formação intelectual. \textbf{Aqui} seria possível observar \textbf{que} o protagonista do conhecimento era o professor, responsável por transmitir os conhecimentos acumulados pela \textbf{humanidade} segundo \textbf{uma} gradação lógica, cabendo aos alunos \textbf{assimilar} os conteúdos \textbf{que} lhes são transmitidos.
Nesse contexto, a prática era determinada pela \textbf{teoria} \textbf{que} a moldava, fornecendo -lhe tanto o conteúdo como a forma de transmissão pelo professor.
No \textbf{que} se refere às correntes educacionais ligadas à renovação \textbf{metodológica}, temos como referenciais pensadores como Rousseau, Kierkegaard, Stirner, Nietzsche e Bergson, chegando ao movimento da Escola Nova, bem como às \textbf{pedagogias} não diretivas, à \textbf{pedagogia} institucional e ao construtivismo priorizando em suas \textbf{teorias} da aprendizagem \textbf{uma} radical mudança, no qual a \textbf{centralidade} é o \textbf{educando}.
Este como protagonista da aprendizagem \textbf{que}, por meio da interação entre si e o professor, constroem seus conhecimentos.
Com esta defesa de \textbf{um} \textbf{método} voltado mais aos educandos, ao professor caberia o papel de acompanha -los de modo a auxiliá -los em seu próprio processo de aprendizagem.
Nessa perspectiva de aprendizagem, o primado é sobre a prática, o eixo do trabalho pedagógico desloca -se da compreensão intelectual para a atividade prática ; dos conteúdos cognitivos para os \textbf{métodos} ou processos de aprendizagem ; do professor para o aluno ; do esforço para o interesse ; da quantidade para a qualidade ( SAVIANI, 2006 ).
Essa tendência torna -se \textbf{hegemônica} sob a forma do movimento da Escola Nova, na segunda metade do \textbf{século} XX, e assume novas versões, entre as quais o construtivismo, sendo a mais difundida na atualidade.
A partir do \textbf{século} XX, a ênfase do processo de aprendizagem se desloca para os \textbf{métodos}, estabelecendo o primado dos \textbf{fundamentos} psicológicos da educação.
Nesse, o professor conduz a aprendizagem, mas o conhecimento tem como ponto de partida a experiência já existente ou a ser realizada pelo próprio aluno, conforme a perspectiva de Piaget ( 1983 ).
Deste modo, pode -se entender \textbf{que} a aprendizagem, no contexto \textbf{teórico} construtivista, está subordinada ao desenvolvimento, ou seja, é sempre provocada por \textbf{uma} situação e \textbf{depende} do desenvolvimento intelectual e da estrutura da própria inteligência.
Na perspectiva construtivista de Piaget, o conhecimento humano se constrói na interação homem-meio, sujeito-objeto.
O aprendizado \textbf{depende} essencialmente de \textbf{uma} prática pedagógica adequada, \textbf{que} favorecesse essa interação, de forma a \textbf{que} gradativamente o aluno pudesse internalizar esse conhecimento, isto é, assimilá -lo integrando -o aos seus demais conhecimentos ( PIAGET, 1975 ).
Conhecer consiste em \textbf{operar} sobre o real e transformá -lo a fim de compreendê -lo, é algo \textbf{que} se dá a partir da ação do \textbf{sujeito} sobre o objeto de conhecimento.
As formas de conhecer são construídas nas trocas com os objetos, tendo \textbf{uma} melhor organização em momentos sucessivos de adaptação ao objeto.
Essa adaptação ocorre através da organização, sendo \textbf{que} o organismo discrimina entre \textbf{estímulos} e sensações, selecionando aqueles \textbf{que} irá organizar em alguma forma de estrutura.
A adaptação possui dois mecanismos opostos, mas complementares, \textbf{que} garantem o processo de desenvolvimento : a \textbf{assimilação} e a acomodação.
Segundo Piaget, o conhecimento é a equilibração/reequilibração entre \textbf{assimilação} e acomodação, ou seja, entre os indivíduos e os objetos do mundo.
A aprendizagem, conforme defende Vygotsky ( 2001, p. 115 ), “ [ … ] pressupõe \textbf{uma} natureza social específica e \textbf{um} processo através do qual as crianças penetram na vida intelectual daqueles \textbf{que} a cercam ”.
É numa relação intrínseca do \textbf{sujeito} com o meio físico e social, orientada por instrumentos e signos ( entre eles a linguagem ), \textbf{que} se processa o seu desenvolvimento cognitivo.
Nessa perspectiva, o desenvolvimento psíquico do \textbf{homem} se realiza por meio de processo de internalização ( VYGOTSKY, 2001 ).
Segundo esse \textbf{autor}, as relações intrapsíquicas ( atividade individual ) constituem -se a partir das relações interpsíquicas ( atividade coletiva ).
É neste movimento do social ao individual \textbf{que} se dá a apropriação de conceitos e significações, ou seja, dá -se a apropriação da experiência social da \textbf{humanidade}.
Desta forma, podemos entender \textbf{que} a aprendizagem não ocorre espontaneamente, ela é mediada culturalmente.
Nas palavras de Leontiev ( 1978, p. 264 ) “ O \textbf{homem} não está evidentemente subtraído ao campo da ação das leis biológicas.
O \textbf{que} é verdade é \textbf{que} as modificações biológicas hereditárias não determinam o desenvolvimento \textbf{sócio-histórico} do \textbf{homem} e da \textbf{humanidade} ”.
A aprendizagem se dá por meio da interação e construção do conhecimento e também dos laços afetivos entre pessoas.
Wallon \textbf{argumenta} \textbf{que} as trocas relacionais da criança com os outros são fundamentais para o desenvolvimento da pessoa.
O meio é \textbf{um} complemento indispensável ao ser vivo.
Ele deverá corresponder às suas necessidades e as suas aptidões sensório-motoras, depois psicomotoras.
Não é mesmo verdadeiro \textbf{que} a sociedade coloca o \textbf{homem} em presença de novos meios, novas necessidades e novos recursos \textbf{que} aumentam a possibilidade de evolução e diferenciação individual.
A constituição biológica da criança ao nascer não será a única do seu destino ( … ).
Os meios em \textbf{que} \textbf{vive} a criança e aqueles com \textbf{que} ela sonha constituem a forma \textbf{que} amolda sua pessoa ( … ) ( WALLON, 1975, p. 164 ).
Para Emília Ferreiro ( 1983 ), a construção do conhecimento tem \textbf{uma} lógica individual na escola e fora dela.
A aprendizagem não é provocada pela escola, mas pela própria \textbf{mente} das crianças, em \textbf{que} esse processo deve ser gradual \textbf{dependendo} de sua \textbf{assimilação} e de \textbf{uma} reacomodação dos esquemas internos, o \textbf{que} \textbf{demanda} tempo.
O processo inicial da aprendizagem é explicado também por variáveis sociais, culturais, políticas e psicolinguísticas.
Complementa a autora, \textbf{que} não existe ponto zero da aprendizagem escrita ; a criança sempre apresenta \textbf{um} conhecimento prévio \textbf{que} o \textbf{sujeito} reestrutura a partir de \textbf{um} processo de acomodação e \textbf{assimilação} mental.
Segundo Antoni Zabala ( 1998 ), a finalidade da escola é promover a formação integral dos alunos.
Para ele, é na instituição escolar, através das relações construídas a partir das experiências \textbf{vividas}, \textbf{que} se estabelecem os vínculos e as condições \textbf{que} definem as concepções pessoais sobre si e os demais.
Daí a necessidade de \textbf{uma} reflexão profunda e permanente da condição de cidadania dos alunos, e da sociedade em \textbf{que} \textbf{vivem}.
Zabala \textbf{atenta} ainda para as particularidades dos processos de aprendizagem de cada aluno.
Na visão de Libâneo ( 1994, p. 88 ), > [ … ] aprender é o processo de \textbf{assimilação} de qualquer forma de conhecimento, daquilo bem \textbf{simples} onde a criança aprende a manipular os brinquedos, aprende a fazer contas, lidar com as coisas, nadar, andar de bicicleta etc., até aquilo mais complexo onde \textbf{uma} pessoa aprende a escolher \textbf{uma} profissão, lidar com as outras.
Dessa forma as pessoas estão sempre aprendendo.
Na origem da aprendizagem está todo \textbf{um} processo de \textbf{assimilação} no qual o aluno, com a orientação do professor, passa a compreender, refletir e aplicar os conhecimentos adquiridos, assim a aprendizagem é observada com a colocação em prática, por parte do estudante, desses conhecimentos \textbf{que} foram transmitidos durante \textbf{uma} aula ou atividade.
A aprendizagem é algo \textbf{que} modifica o pensamento.
Para \textbf{que} se possa haver a aprendizagem o aluno \textbf{necessita} ser estimulado com conteúdo de seu alcance, textos \textbf{que} tratem de sua realidade.
Nas palavras de Libâneo ( 1994, p. 90 ) “ a relação entre ensino e aprendizagem não é algo mecanizado, não é apenas \textbf{uma} transmissão feita pelo professor \textbf{que} ensina para \textbf{um} aluno \textbf{que} aprende ”.
Tem para ele, ao contrário, \textbf{um} dever de “ estabelecer exigências e expectativas \textbf{que} os alunos possam cumprir e, com isso, mobilizem suas energias.
Esse sim, pois é o seu papel, impulsionar a aprendizagem, precedendo -a muitas das vezes. ” Segundo Paulo Freire ( 1996, p. 14 ), não existe ensino sem aprendizagem, e nesse processo, “ educador e educandos, \textbf{lado} a \textbf{lado}, vão se transformando em reais sujeitos da (re)construção do saber, pois o conhecimento não está no professor, o conhecimento circula, é compartilhado ”.
O processo de ensinar tem por finalidade o alcance da aprendizagem pelo estudante.
Nesse sentido, na condução satisfatória do processo de ensino, efetivo e \textbf{que} \textbf{agregue} valor a vida dos estudantes, o professor adquire o papel de mediador e facilitador do processo de \textbf{ensino-aprendizagem}, assim precisa utilizar \textbf{métodos} adequados e centrados na realidade dos estudantes, num diálogo constante entre estes \textbf{atores}, buscando estabelecer e fortalecer vínculos.
Esse estreitamento de vínculos deve ser ponderado por regras onde todos têm liberdade, \textbf{uma} vez \textbf{que}, como afirma Freire ( 1986. p.54 ) : “ \textbf{uma} experiência \textbf{dialógica} \textbf{que} não se \textbf{baseia} na seriedade e na competência é muito pior do \textbf{que} \textbf{uma} experiência ‘ bancária ’, onde o professor simplesmente transfere conhecimento ”.
Com isto, pode ser afirmado \textbf{que} nesse diálogo, o encontro das pessoas ocorre para \textbf{que} elas reflitam sobre o mundo, sobre como transformá -lo e melhorá -lo, tendo na sala de aula o seu primeiro lugar de exercício prático.
Frente a isto, objetiva -se \textbf{uma} educação \textbf{que} ultrapasse a \textbf{mera} formação \textbf{conceitual}, conduzindo o \textbf{alunado} à reflexão sobre seu papel na sociedade, para \textbf{que} ele possa nessa relação \textbf{teoria} e prática encontrar novas maneiras de atuar no contexto em \textbf{que} \textbf{vive}. \textbf{Uma} Educação \textbf{que} tenha a intenção de formar no aluno \textbf{uma} atitude crítica quanto às relações sociais nas quais se insere.
Constatar essas relações possibilita colocar “ nas mãos dos educadores \textbf{uma} arma de luta capaz de permitir -lhes o exercício de \textbf{um} poder real, ainda \textbf{que} limitado ” ( SAVIANI, 2003, p.31 ), \textbf{uma} vez \textbf{que} a educação se torna decisiva neste constante processo de transformação da realidade social.
A criação de sentido pessoal e social, orientada por princípios de solidariedade e favorecida pelas interações \textbf{que} aumentam o aprendizado instrumental produz a aprendizagem \textbf{dialógica}.
Essa aprendizagem apresenta -se como ferramenta \textbf{que} vai auxiliar na superação de desafios \textbf{que} estão internalizados ao bom êxito do processo de ensino e aprendizagem, \textbf{visto} \textbf{que}, essa aprendizagem fundamenta -se e expressa -se por meio do diálogo igualitário e na interação mais diversa possível entre as pessoas.
Como na \textbf{contemporaneidade} muito se tem refletido a respeito de como deve ser a atuação do professor em sala de aula no processo de ensino/aprendizagem, reafirmamos \textbf{que} cabe ao professor, na coordenação do processo de aprendizagem, articular os conhecimentos científicos, tecnológicos e artísticos, cientes dos vários elementos didáticos \textbf{que} \textbf{influenciam} e possibilitam a apropriação dos conhecimentos pelos alunos, seguindo \textbf{uma} perspectiva \textbf{dialética} entre o saber \textbf{imediato} e o mediato.
Com adoção de prática pedagógica \textbf{que} tenha dinâmica própria, \textbf{que} permita o exercício do pensamento reflexivo, conduza a \textbf{uma} visão política de cidadania e \textbf{que} seja capaz de integrar a arte, a cultura, os valores e a interação, propiciando, assim, a autonomia dos alunos e a reflexão sobre seu lugar no mundo, de forma significativa.

% matched lemmas: agregar, alunado, antropológico, aproximar, aqui, argumentar, assimilar, assimilação, atentar, ator, autor, basear, centralidade, conceitual, conhecido, contemporaneidade, demandar, depender, dialético, dialógico, dizer, educar, embasar, enquanto, ensino-aprendizagem, estímulo, exposição, focar, fundamento, hegemônico, homem, humanidade, imediato, influenciar, justo, lado, mente, mero, metodológico, moderno, método, necessitar, operar, pedagogia, que, revelar, simples, sujeito, século, sócio-histórico, teoria, teórico, um, ver, vertente, viver
\end{textsample}
